\section*{第 \ref{ch:g12:diffeq} 章 --- 微分方程}

\begin{solution}{exo:g12:diffeq:1}
\emph{(a)} 由 $y = C\eu^{3x}$；$y(0) = 2$ 得 $y = 2\eu^{3x}$。

\emph{(b)} $y' = -\frac12 y$，于是 $y = C\eu^{-x/2}$；$y(0) = -1$ 得
$y = -\eu^{-x/2}$。

\emph{(c)} 平衡解 $y_p = 5$；$y = C\eu^{-x} + 5$；$y(0) = 0$ 得
$C = -5$：$y = 5\left(1 - \eu^{-x}\right)$。
\end{solution}

\begin{solution}{exo:g12:diffeq:2}
$N' = kN$，故 $N(t) = N_0 \eu^{kt}$。经过 $3$ 小时数量加倍，即
$\eu^{3k} = 2$，\emph{即}
\[
k = \frac{\ln 2}{3} \approx 0.231\ \text{h}^{-1}.
\]
\end{solution}

\begin{solution}{exo:g12:diffeq:3}
$y_p'(x) = \eu^x + x\eu^x = y_p(x) + \eu^x$：$y_p$ 是一个特解
（这正是 \cref{met:g12:diffeq:particular} 所述的共振情形）。由
\cref{prop:g12:diffeq:general}，方程的解为
$y = C\eu^{x} + x\eu^{x} = (C + x)\,\eu^x$，$C \in \R$。
\end{solution}

\begin{solution}{exo:g12:diffeq:4}
试取 $y_p = \beta\eu^x$：$\beta\eu^x = -2\beta\eu^x + \eu^x$ 得
$3\beta = 1$，于是 $y_p = \frac13\eu^x$。通解
为 $y = C\eu^{-2x} + \frac13\eu^{x}$；由条件 $y(0) = 1$ 得
$C = \frac23$：
\[
y(x) = \frac{2}{3}\,\eu^{-2x} + \frac{1}{3}\,\eu^{x}.
\]
\end{solution}

\begin{solution}{exo:g12:diffeq:5}
$N(t) = N_0\,\eu^{-\lambda t}$，其中 $\lambda = \frac{\ln 2}{5730}$
（见\cref{ex:g12:exp:halflife}）。解方程 $\eu^{-\lambda t} = 0.6$：
\[
t = \frac{\ln(1/0.6)}{\lambda} = 5730\,\frac{\ln(5/3)}{\ln 2}
\approx 5730 \times \frac{0.5108}{0.6931} \approx 4220 \text{ 年}.
\]
\end{solution}

\begin{solution}{exo:g12:diffeq:6}
\emph{1.} 盐的流入速率为 $0.2 \times 5 = 1$ kg/min。流出液
的浓度为 $\frac{m}{100}$ kg/L，流出速率为 $5$ L/min，\emph{即} $\frac{m}{20}$
kg/min。于是 $m' = 1 - \frac{m}{20}$。

\emph{2.} 平衡解 $m_p = 20$；$m(t) = 20 + C\eu^{-t/20}$，且
$m(0) = 0$ 得 $C = -20$：
\[
m(t) = 20\left(1 - \eu^{-t/20}\right) \xrightarrow[t\to+\infty]{} 20 \text{ kg}.
\]
长期来看，水箱中的浓度与流入盐水的浓度相等：
$0.2$ kg/L $\times$ $100$ L $= 20$ kg。
\end{solution}

\begin{solution}{exo:g12:diffeq:7}
\emph{1.} $v' = g - \frac{k}{m} v$，其中 $\frac km = \frac{16}{80} = 0.2$。
平衡解 $v_\infty = \frac{mg}{k} = \frac{9.8}{0.2} = 49$ m/s；
$v(t) = 49 + C\eu^{-0.2t}$，由 $v(0) = 0$ 得
\[
v(t) = 49\left(1 - \eu^{-0.2 t}\right).
\]

\emph{2.} 终速度为 $49$ m/s（$\approx 176$ km/h）。要使
$1 - \eu^{-0.2t} = 0.95$，\emph{即} $\eu^{-0.2t} = 0.05$：
\[
t = \frac{\ln 20}{0.2} \approx \frac{3.00}{0.2} \approx 15 \text{ s}.
\]
\end{solution}

\begin{solution}{exo:g12:diffeq:8}
\emph{1.} 令 $z = \frac1y$，则 $z' = -\frac{y'}{y^2}
= -\frac{y(1-y)}{y^2} = -\frac{1-y}{y} = -\frac1y + 1 = -z + 1$。

\emph{2.} 平衡解 $z_p = 1$，故 $z(t) = 1 + C\eu^{-t}$。由
$y(0) = \frac{1}{10}$，$z(0) = 10$，于是 $C = 9$，且
\[
y(t) = \frac{1}{1 + 9\,\eu^{-t}} .
\]

\emph{3.} 当 $t \to +\infty$ 时，$9\eu^{-t} \to 0$，$y(t) \to 1$。
曲线为经典的 S 形（\emph{sigmoid}）逻辑斯谛曲线：起初增长缓慢
（$y$ 较小，$y' \approx y$），增长最快出现在 $y = \frac12$ 时
（此时 $y' = y(1-y)$ 最大），随后逐渐饱和，
趋向环境容纳量 $1$。
\end{solution}

\begin{solution}{exo:g12:diffeq:9}
由\cref{thm:g12:diffeq:affine}，$y(x) = C\eu^{ax} - \frac ba$。于是
\[
u_{n+1} = C\eu^{a(n+1)} - \frac ba
= \eu^{a}\left(C\eu^{an} - \frac ba\right) + \frac ba\left(\eu^a - 1\right)
= q\,u_n + r
\]
其中 $q = \eu^{a}$，$r = \frac{b}{a}\left(\eu^{a} - 1\right)$。
条件 $\abs q < 1$ 等价于 $\eu^a < 1$，\emph{即} $a < 0$：这恰是使
微分方程的解收敛于平衡解
$-\frac ba$ 的条件 --- 事实上，递推关系的不动点
为 $\frac{r}{1 - q} = -\frac{b}{a}$。
\end{solution}

\begin{solution}{pb:g12:diffeq:1}
\textbf{1.} $y = 2\eu^{3t}$。第二小题的平衡解为 $3$，
故 $y = 3 + C\eu^{-2t}$，由 $y(0) = 0$ 得 $C = -3$：
$y = 3\left(1 - \eu^{-2t}\right)$。

\textbf{2.} $\left(y\eu^{-at}\right)' = y'\eu^{-at} -
ay\eu^{-at} = (y' - ay)\eu^{-at} = 0$：乘积为
常数 $C$，故 $y = C\eu^{at}$ --- 没有解能逸出这一形式。

\textbf{3.} 平衡解：$y \equiv 3$。其余解均为
$3 + C\eu^{-2t}$，其中 $C \neq 0$：指数项消失，解
平滑地趋于 $3$ --- 一个稳定不动点，
这是递推不动点在连续情形下的对应，见
\cref{pb:g11:seq:1}。

\textbf{4.} 当 $y = y_0\eu^{-t/\tau}$ 减半时，
$\eu^{-t/\tau} = \frac12$：$t = \tau\ln 2$。

\textbf{5.} 所有解都是衰减曲线的竖直平移，
它们都朝 $3$ 衰减：起点在上方者降至 $3$，起点在下方者
升至 $3$，常数解 $3$ 静止不动 --- 一个
围绕平衡点的漏斗。

\textbf{6.} $\lambda = \frac{\ln 2}{5730} \approx 1.21 \times
10^{-4}$ 每年。

\textbf{7.} $\eu^{-\lambda t} = 0.2$:
$t = \frac{\ln 5}{\lambda} \approx 13\,300$ 年。

\textbf{8.} $t = \frac{\ln(1/0.15)}{\lambda} \approx 15\,700$
年：拉斯科洞穴的公牛壁画属于冰河时代晚期。

\textbf{9.} $t = \frac{\ln(1/0.53)}{\lambda} \approx 5\,250$
年：与考古学家的估计值相差在一个人寿命之内 ---
时钟可靠。

\textbf{10.} 十个半衰期后剩下 $2^{-10} \approx 0.1\,\%$：
处于测量能力的边缘。经过 $65$ 百万年后，
剩余比例为 $2^{-65\,000\,000/5\,730} \approx 2^{-11\,344}$：
任何化石中都不再残留原始存量中的原子 ---
恐龙需用更慢的时钟定年（钾--氩法及
同类方法）。

\textbf{11.} 取 $52\,\%$ 时得 $\approx 5\,405$ 年，
取 $54\,\%$ 时得 $\approx 5\,095$：化学中的 $\pm 1\,\%$
在历史年代上约合 $\pm 155$ 年 --- 实验室中的精度
就是博物馆标签上的精度。

\textbf{12.} $z = T - T_a$ 满足 $z' = -kz$：
$z = (T_0 - T_a)\eu^{-kt}$，且 $T = T_a + z$。

\textbf{13.} 初始温差为 $T_0 - T_a = 70$；五
分钟后温差为 $70 - 20 = 50$，故 $50 = 70\eu^{-5k}$：
$k = \frac{\ln(70/50)}{5} \approx 0.067$ 每分钟。可饮用时：
$55 - 20 = 35 = 70\eu^{-kt}$:
$t = \frac{\ln 2}{k} \approx 10.3$ 分钟。

\textbf{14.} 散热速率与温差
$T - T_a$ 成正比：刚冲泡的滚烫咖啡散热最快。
\emph{此时}加入牛奶，温差立即骤减，故
十分钟内散失的热量较少；若最后才加，
咖啡先全速降温。要喝到最热的一杯，应
先加牛奶。（性急的主人搞反了热力学。）

\textbf{15.} 与环境的温差：午夜时为 $10$，一小时后
为 $8$：$\eu^{-k} = 0.8$：
$k = \ln\frac{10}{8} \approx 0.223$ 每小时。死亡时的温差
为 $17$；至午夜时已降为 $10$：
$s = \frac{\ln(17/10)}{k} \approx 2.4$ 小时。死亡时刻：
约 $21{:}40$。

\textbf{16.} 房间若有供暖或穿堂风（$T_a$ 非常数），
两种情况都会使时钟走偏；衣物或体重会改变 $k$
（法医的查表会对此加以修正） --- 错误的 $k$ 会使
整个时间区间按比例失真；被从别处挪来的尸体会
在降温中途重设 $T_a$，使死亡时间显得或早或晚。
方程本身是可靠的；风险在于假设。

\textbf{17.} $z = \frac1y$ 满足 $z' = 1 - z$：
$z = 1 + 9\eu^{-t}$（由 $z(0) = 10$ 得），故
$y = \frac{1}{1 + 9\eu^{-t}}$。拐点位于 $y = \frac12$ 处：
$9\eu^{-t} = 1$，即 $t = \ln 9 \approx 2.2$ --- 疫情的
转折日期，直接由公式得出。

\textbf{18.} 终端速度：$v' = 0$ 时 $v = 20$ m/s。
解为 $v(t) = 20\left(1 - \eu^{-t/2}\right)$。于是
达到终端速度的 $0.95$ 倍时，$\eu^{-t/2} = 0.05$：$t = 2\ln 20 \approx 6$ s ---
雨滴在数秒之内即达到终端速度，这正是
雨水不致命的原因。

\textbf{19.} 稳态：$c = 8$。其一半为
$4 = 8\left(1 - \eu^{-t/2}\right)$，即 $\eu^{-t/2} = \frac12$：
$t = 2\ln 2 \approx 1.4$ 小时。点滴问题是
贷款问题的连续版本：流入恒定，流出与当前量成正比 ---
\cref{exo:g12:diffeq:9} 使这套对应关系严格成立。

\textbf{20.} 衰减（$a < 0$，$b = 0$）、冷却（$y = T - T_a$），
下降（$a < 0$，$b > 0$：从
下方趋近平衡）、静脉给药（同上，经静脉）：一个方程，四种情境。
循环步骤：陈述变化规律；列出方程；求解
（\cref{met:g12:diffeq:solve}）；依测量数据标定
常数；作出预测；然后检验假设（第
16 题）。振荡需满足 $y'' = -\omega^2 y$ --- 此方程已遇到、已求解，
并在\cref{pb:g12:trigo:1}中使之振动起来。
\end{solution}
